% Methods: environment and reproducibility, the generated platform, the 10-gate
% design, the wired standards corpus, the pipeline, the iteration design, and the
% test floor. Built by the future pass from the instructions below. Do not process
% now. Figure 05 (standards wheel) is placed here; \autoref it from the prose.

[Write the Methods as a connected narrative across subsections, each carrying a
left-aligned fixed-width table for the dense facts so the prose stays readable.
Every table column must use the format \texttt{\{>\{\textbackslash raggedright
\textbackslash arraybackslash\}p\{width\}\}}. Where a table lists generated files,
put the common parent path once in the caption or the top row and abbreviate the
leaf names in the cells. Reference the tables and the figure; do not narrate them
row by row.]

[Subsection 1, Environment and reproducibility. Source this verbatim in substance
from \texttt{papers/VVUQ-02/execution/01-foundation/environment-and-verification.md}.
Report the execution host (Linux 6.18.5 managed ephemeral cloud container), the
interpreter (CPython 3.11.15), the autonomous runner (Claude Code Opus 4.7 1M Max),
and the dependency posture (numpy, click, pyyaml, jsonschema core plus ruff,
yamllint, pytest), with the heavy optional groups for anthropic, ollama, openai,
and live Zenodo deliberately not installed so the guarded offline paths ran. State
the determinism result as the credibility floor: the sensor CSV, the comparison,
and the sweep index each regenerated byte-for-byte or as an identical semantic
match from seed 20260525, which is exactly what NASA-STD-7009A and ASME V\&V 40-2018
§3 expect before any verification or validation claim is credible. Cite
\texttt{nasastd7009a} and \texttt{asmevv40}. Build one table of the regenerated
artifacts and their commands from the determinism table in that source file.]

[Subsection 2, The generated platform. Describe what the codegen tree models, from
\texttt{papers/VVUQ-02/codegen/README.md} and the platform docs in
\texttt{papers/VVUQ-02/codegen/docs/} (h2\_robot\_specification.md,
h2\_hand\_kinematics.md, sensor\_spec.md, perception\_stack\_spec.md,
suturing\_anastomosis\_protocol.md): the hypothetical 2030 Unitree H2-Surgical 1.0,
71 total degrees of freedom, performing a 60-second eight-phase Whipple. Summarise
the source package map from the \texttt{Source executed} table in
\texttt{papers/VVUQ-02/execution/README.md} (the \texttt{src/} subpackages
kinematics, sensors, perception, autonomy, hands, balance, safety, suturing, vvuq,
simulation, metrics, llm, zenodo) and the configuration set in
\texttt{papers/VVUQ-02/codegen/config/} (h2\_kinematics.yaml, hand\_model.yaml,
balance\_model.yaml, perception\_model.yaml, autonomy\_plan.yaml,
safety\_zones.yaml, shared\_or\_actors.yaml, anastomosis\_targets.yaml,
project.yaml). Build one platform-summary table; abbreviate the config leaf names
and give the \texttt{papers/VVUQ-02/codegen/config/} parent path once at the top.]

[Subsection 3, The ten-gate assurance design. This is the methodological core.
From \texttt{papers/VVUQ-02/codegen/config/vvuq\_thresholds.yaml},
\texttt{papers/VVUQ-02/codegen/docs/vvuq\_gate\_spec.md}, and the threshold table in
\texttt{papers/VVUQ-02/execution/03-vvuq/README.md}, explain the Verify then
Validate then Quantify then Evaluate composition implemented in
\texttt{papers/VVUQ-02/codegen/src/vvuq/vvuq\_gate.py} and
\texttt{gate\_registry.py}: a verification fraction that must equal 1.0, a
validation agreement against an independent reference in
\texttt{papers/VVUQ-02/codegen/data/reference/}, a coefficient-of-variation bound
across three or more seeded runs, and an extra hard predicate on the three
immediate-catastrophe gates (06 vascular no-fly, 09 shared-OR collision, 10 fault
e-stop). Explain that bounds are set in proportion to model risk per ASME V\&V 40,
and cite \texttt{asmevv40} and \texttt{iso14971}. Do not reproduce all ten
threshold rows in the prose; that table belongs to Results alongside the forest
plot, and should be cross-referenced here.]

[Subsection 4, The wired external-standards corpus. Explain that credibility comes
from binding each gate to published consensus standards already used in real life,
resolved at runtime, not asserted. Source this from
\texttt{papers/VVUQ-02/inputs/README.md},
\texttt{papers/VVUQ-02/inputs/standards/manifest.yaml}, the per-standard summaries
in \texttt{papers/VVUQ-02/inputs/standards/} (asme\_vv40\_2018.md,
nasa\_std\_7009a.md, iec\_80601\_2\_77\_2019.md, iso\_13482\_2014.md,
iso\_ts\_15066\_2016.md), the clinical baselines in
\texttt{papers/VVUQ-02/inputs/clinical/} (dutch\_cohort\_2025.md,
fistula\_risk\_score.md), and the machine-readable binding in
\texttt{papers/VVUQ-02/codegen/config/standards\_map.yaml}. Use this subsection to
introduce \autoref{fig:wheel}, which organises the corpus into six domains around
the assurance hub. Cite the standards being introduced: \texttt{asmevv40},
\texttt{nasastd7009a}, \texttt{iec8060127}, \texttt{iec60601},
\texttt{isots15066}, \texttt{iso13482}, \texttt{iso10218}, \texttt{iso9283},
\texttt{iec62304}, \texttt{iso14971}, \texttt{iso13849}, \texttt{ul4600},
\texttt{ieee7009}, and \texttt{fda2023credibility}.]

[Subsection 5, The control pipeline and the iteration design. From
\texttt{papers/VVUQ-02/execution/02-pipeline/README.md} and its artifact
\texttt{papers/VVUQ-02/execution/02-pipeline/artifacts/pipeline\_execution\_log.txt},
describe the intent to compile to act to score flow across the six behavior groups
(autonomy, kinematics, perception, hands, balance, suturing), noting that the
on-prem LLM intent ran on its deterministic offline fallback. From
\texttt{papers/VVUQ-02/codegen/config/project.yaml} and
\texttt{papers/VVUQ-02/codegen/src/simulation/iterate.py}, state the 32-iteration
Latin hypercube at seed 20260525 over five free parameters, and the six-component
composite in \texttt{papers/VVUQ-02/codegen/src/metrics/compute.py}. Reproduce the
\texttt{External standards anchored here} mapping from the pipeline README as a
table binding each behavior group to its gate and governing standards.]

[Subsection 6, The verification floor. From
\texttt{papers/VVUQ-02/execution/01-foundation/test-suite.md} and the test modules
in \texttt{papers/VVUQ-02/codegen/tests/}, state that the shipped suite is 172
tests across 15 modules, that the assurance suite \texttt{test\_vvuq\_gates.py}
alone carries 64 tests (more than a third of the budget, itself evidence for the
thesis), and that the suite maps to ASME V\&V 40-2018 §8 verification and IEC 62304
software verification. Cite \texttt{iec62304}. Build the per-module test table from
that source, abbreviating module names. Note that local execution used a single
interpreter and that cross-version coverage is delegated to the CI matrix, which
belongs to Limitations.]

\begin{figure}[H]
   \centering
   \IfFileExists{Images/05-clinical-regulatory-standards-wheel.png}{%
     \hspace{-0.95cm}%
     \includegraphics[width=1.1\linewidth, trim={0 0 0 0}, clip]{Images/05-clinical-regulatory-standards-wheel.png}%
     % \vspace{-0.25cm}
   }{%
     \setlength{\fboxsep}{10pt}%
     \fbox{\parbox[c][6cm][c]{0.9\linewidth}{\centering\itshape Figure placeholder.
     The author adds \texttt{Images/05-clinical-regulatory-standards-wheel.png},
     rendered by
     \texttt{papers/VVUQ-02/imagegen/05-clinical-regulatory-standards-wheel/05-clinical-regulatory-standards-wheel.py}.}}%
   }
    \caption{Clinical and regulatory standards wheel for the wired VVUQ-02 input corpus of fourteen consensus standards and two clinical baselines across six domains.}
    \label{fig:wheel}
\end{figure}
